%% Approximation du signal carre par sa serie de Fourier tronquee a
%% N=1 (pointille gris), N=5, N=9 harmoniques -- convergence + phenomene
%% de Gibbs visible pres des sauts. Le vrai signal carre (noir) est
%% trace en reference, cible vers laquelle les troncatures convergent.
%% Consomme par slides.
\begin{tikzpicture}
    \begin{axis}[
        width=10cm, height=5.5cm,
        ytick=\empty, xtick=\empty,
        axis lines=middle,
        ymin=-1.5, ymax=1.5,
        domain=0:4*pi,
        samples=200,
    ]
    \addplot[black, thick] coordinates {
        (0,1) (3.14159,1) (3.14159,-1) (6.28319,-1)
        (6.28319,1) (9.42478,1) (9.42478,-1) (12.56637,-1)
    };
    \addplot[gray!70, dashed, thick] {(4/pi)*sin(deg(x))};
    \addplot[orange, thick] {(4/pi)*(sin(deg(x)) + (1/3)*sin(deg(3*x)) + (1/5)*sin(deg(5*x)))};
    \addplot[niceblue, very thick] {(4/pi)*(sin(deg(x)) + (1/3)*sin(deg(3*x)) + (1/5)*sin(deg(5*x)) + (1/7)*sin(deg(7*x)) + (1/9)*sin(deg(9*x)))};
    \end{axis}
\end{tikzpicture}
